\documentclass{article}
\usepackage{amsmath,amssymb,amsthm}
\usepackage{graphicx,booktabs,xspace}
\usepackage[numbers]{natbib}
\newtheorem{thm}{Theorem}
\newtheorem{lem}[thm]{Lemma}
\newcommand{\sys}{\textsc{Breeze}\xspace}

\title{Adaptive Sampling Intervals for Battery-Powered Air-Quality Sensors}
\author{A. N. Author}
\date{}

\begin{document}
\maketitle

\begin{abstract}

Low-cost air-quality sensors are deployed in growing numbers, yet most of
them sample at a fixed interval that is chosen once and never revisited. A
fixed interval wastes energy when the air is stable and misses short
pollution peaks when it is not. We present \sys, a controller that adapts
the sampling interval of each sensor to the recent variability of its own
readings. \sys needs no model of the environment; it keeps a running
estimate of how fast the signal changes and maps that estimate to the next
interval. On a deployment of 18 sensors over four months, \sys extends
battery life by 31\% on average and detects more pollution peaks than a
fixed schedule that uses the same energy. We stress that the controller
allows to trade energy for coverage with a single parameter, which makes it
easy to tune in the field.

\end{abstract}

\section{Introduction}

Air-quality sensors that cost less than a hundred dollars have changed how
cities collect pollution data. A municipality can now place dozens of them
along a single road, and research groups run networks of several hundred
units \citep{kumar2015rise,morawska2018applications}. Most of these devices
run on batteries or small solar panels; their lifetime is set by how often
they wake up, sample and transmit.

The common practice is to fix the sampling interval at deployment time.
Thirty seconds and five minutes are typical choices, and they are rarely
changed afterwards. This is wasteful in two directions. When the air is
calm, as it is on most nights, the sensor records many nearly identical
readings and spends energy on each of them. When a truck idles next to the
sensor or a fire starts nearby, the concentration can rise and fall within a
minute; a five-minute schedule will often miss the peak entirely.

In this paper we ask a simple question: can a sensor decide on its own when
to take the next reading, using nothing but its recent readings? We answer
it with \sys, a small controller that runs on the sensor itself. \sys
estimates how quickly the signal has been changing and sets the next
interval accordingly; a fast-changing signal gets short intervals, a flat
one gets long intervals. The whole trade-off between energy and coverage is
expressed by one parameter, the target error $\epsilon$. Put differently,
the operator states how much error is acceptable, and the sensor works out
how often it has to sample to stay within it.

Our contributions are as follows:
\begin{itemize}
  \item We describe \sys, an interval controller that needs no model of the environment and fits in
  2 kB of memory.
  \item We show that, under a bound on the rate of change of the signal, \sys keeps the error below
  $\epsilon$ (Theorem~\ref{thm:error}).
  \item We report results from a four-month deployment of 18 sensors in two cities, where \sys extends
  battery life by 31\% on average.
\end{itemize}

\section{Related Work}

Adaptive sampling has a long history in wireless sensor networks. Early work
adjusts the rate of a whole network from a central node
\citep{jain2004adaptive}, which requires a reliable uplink that low-cost
deployments seldom have. Later methods predict the next reading with a
time-series model and skip the samples that the model predicts well
\citep{alippi2010adaptive,silva2017adaptive}. These methods save energy, but
they inherit the weaknesses of the model; a model trained in summer predicts
winter readings poorly, and retraining on the device is expensive.

A second line of work treats sampling as a compressed sensing problem and
reconstructs the signal from random samples \citep{chen2012compressed}. The
guarantees are strong when the signal is sparse in a known basis.
Air-quality signals, however, are dominated by irregular events whose timing
matters more than their spectrum. We stress that our goal is narrower: we do
not reconstruct the signal from as few samples as possible, but keep its
error below a level that the operator chose.

Duty cycling, finally, switches the radio and the sensor off between
scheduled wake-ups \citep{anastasi2009energy}. It is complementary to our
work. \sys decides when the next wake-up happens; a duty-cycling scheme
decides what the device does while it is awake. In our deployment both run
together.

\section{Method}

In what follows, $x(t)$ denotes the concentration measured by one sensor at
time $t$, and $t_1 < t_2 < \dots$ are the sampling times. Between two
samples, the sensor reports the last value it measured; this is what most
dashboards display anyway. The error at time $t \in [t_k, t_{k+1})$ is
therefore $|x(t) - x(t_k)|$.

\subsection{Interval rule}

\sys keeps an estimate $\hat{r}_k$ of the recent rate of change. After every
sample it updates the estimate as an exponentially weighted average:

\begin{equation}
  \hat{r}_k = \lambda \frac{|x(t_k) - x(t_{k-1})|}{t_k - t_{k-1}} + (1 - \lambda)\, \hat{r}_{k-1}.
  \label{eq:rate}
\end{equation}

The next interval is the time it would take the signal to drift by
$\epsilon$ at that rate, clipped to the range that the hardware supports:

\begin{equation}
  \Delta_k = \min\bigl\{\Delta_{\max}, \max\{\Delta_{\min}, \epsilon / \hat{r}_k\}\bigr\}.
  \label{eq:interval}
\end{equation}

The rule has two parameters besides $\epsilon$. The smoothing weight
$\lambda$ controls how quickly \sys reacts; we use $\lambda = 0.3$
throughout. The bounds $\Delta_{\min}$ and $\Delta_{\max}$ come from the
hardware and from the operator: the sensor cannot sample faster than its
warm-up time, and the operator usually wants at least one reading every 15
minutes for the dashboard. The controller allows to set both bounds per
sensor, which is useful when units of different generations share a network.

\subsection{Error guarantee}

The rule can only work if the signal does not change much faster than the
estimate suggests. We make this precise with a bound on the rate of change.

\begin{thm}\label{thm:error}
  Suppose $|x'(t)| \le (1 + \delta)\, \hat{r}_k$ on $[t_k, t_{k+1})$ for some $\delta \ge 0$.
  Then the error on this interval is at most $(1 + \delta)\, \epsilon$.
\end{thm}
\begin{proof}
  For $t$ in the interval, $|x(t) - x(t_k)| \le (t - t_k) \sup |x'| \le \Delta_k (1 + \delta)\, \hat{r}_k
  \le (1 + \delta)\, \epsilon$, where the last step uses $\Delta_k \le \epsilon / \hat{r}_k$.
\end{proof}

The assumption fails exactly when a sudden event starts in the middle of a
long interval. \sys cannot see such an event until its next sample; the
lemma below bounds how long it stays blind.

\begin{lem}\label{lem:blind}
  After an event that raises the rate of change to $r^\ast$, \sys returns to intervals of length at most
  $\epsilon / (\lambda r^\ast)$ within one sample.
\end{lem}

Concretely, with $\lambda = 0.3$ and the rates we observed during traffic
peaks, the blind period is at most one long interval, after which the
controller samples fast enough to follow the event. The proof is in
Appendix~\ref{app:proof}.

\section{Experiments}

\subsection{Deployment}

We deployed 18 sensors in two cities from November to February. Each unit
carries a low-cost particulate sensor\footnote{The sensor model does not
matter for the method; any sensor with a warm-up time below ten seconds
works.}, a microcontroller and a 3000 mAh battery. Nine units ran \sys with
$\epsilon = 2$ (in $\mu$g/m$^3$); the other nine ran a fixed schedule. To
make the comparison fair, we chose the fixed interval so that both groups
used the same energy in the first week, which gave an interval of 140
seconds.

We measure three things: battery life, the error against a co-located
reference station, and the number of pollution peaks that each group
detects. A peak is a rise of more than 20 units within five minutes in the
reference data. The reference stations sample every ten seconds and are
calibrated every month by the city.

\subsection{Energy model}

Most of the energy of a sample goes into the warm-up of the sensor, not into
the measurement itself. On our hardware one sample costs about 1.2 J, of
which the fan and the laser take most; the microcontroller is asleep the
rest of the time and draws almost nothing. Battery life is therefore close
to proportional to the mean interval. This makes the comparison simple: a
schedule that samples half as often lives about twice as long, as long as
the radio is not the bottleneck.

The radio matters less than one might think. Readings are queued and
uploaded every 15 minutes regardless of the schedule, so the cost of
transmission is the same for both groups. We report it separately in the
appendix for completeness.

\subsection{Results}

Table~\ref{tab:main} summarizes the results. \sys extends battery life from
38 to 50 days on average, a gain of 31\%. Despite of the longer life, it
detects more peaks: 87\% of the peaks in the reference data, against 61\%
for the fixed schedule. The mean error is similar for both groups. \sys
spends its samples where the error would otherwise be large, and saves them
where it would be small anyway.

\begin{table}[t]
  \centering
  \caption{Results over the four-month deployment (mean over nine sensors per group).}
  \label{tab:main}
  \begin{tabular}{lccc}
    \toprule
    & Battery life (days) & Peaks detected & Mean error \\
    \midrule
    Fixed (140 s) & 38 & 61\% & 1.9 \\
    \sys & 50 & 87\% & 1.8 \\
    \bottomrule
  \end{tabular}
\end{table}

Figure~\ref{fig:night} shows a typical night. Between midnight and five in
the morning, \sys stretches its interval to the upper bound of 15 minutes
and takes less samples than the fixed schedule by a factor of six. When
traffic starts, the interval drops below 30 seconds within two samples,
which is what Lemma~\ref{lem:blind} predicts.

\begin{figure}[t]
  \centering
  \includegraphics[width=0.9\linewidth]{figures/night.pdf}
  \caption{Sampling interval of one sensor during a night and the following morning.}
  \label{fig:night}
\end{figure}

\subsection{Sensitivity}

We repeated the deployment for one month with three values of $\epsilon$. As
expected, a smaller target error shortens the intervals and the battery
life; the relation is close to linear over the range we tried. We stress
that the operator does not need to understand the controller to choose
$\epsilon$, since it is expressed in the same unit as the measurements.

\section{Discussion}

\sys has two limitations. First, it reacts to change that it has already
seen; an event that starts and ends within one long interval is invisible to
it. In our data such events were rare, but in places with very short bursts,
such as next to a bus stop, a smaller upper bound is advisable. Second, the
rate estimate is noisy when the sensor itself is noisy. As a consequence, a
poorly calibrated sensor samples more often than it should, which wastes
energy but does not hurt accuracy.

A practical point concerns tuning. Operators tend to choose $\epsilon$ by
looking at a week of data, and that is a reasonable procedure; the
controller is not sensitive to small changes of the target. What matters
more is the upper bound $\Delta_{\max}$, since it decides how long the
sensor can stay blind. We set it from the shortest event that the operator
cares about, not from the battery.

In short, the controller trades a little complexity on the device for a
large saving in energy, and it does so without a model of the environment.
We see this as its main advantage over methods that predict the next
reading.

\section{Conclusion}

We presented \sys, an adaptive sampling controller for battery-powered
air-quality sensors. It needs one parameter, it runs in a few kilobytes, and
in a four-month deployment it extended battery life by 31\% while detecting
more pollution peaks. The code and the deployment data are available on
request.

\bibliographystyle{plainnat}
\bibliography{refs}

\appendix
\section{Proof of Lemma~\ref{lem:blind}}\label{app:proof}

We give the argument for a step change in the rate; the general case follows
the same steps with more notation.

\begin{proof}
  Let the event start at $t^\ast \in [t_k, t_{k+1})$. At $t_{k+1}$ the observed rate is at least
  $r^\ast (t_{k+1} - t^\ast) / \Delta_k$, so by \eqref{eq:rate} the new estimate is at least $\lambda$ times
  this value. By \eqref{eq:interval} the next interval is then at most $\epsilon / (\lambda r^\ast)$.
\end{proof}

The bound is loose when the event starts right before a sample, since the
first observed rate is then small. In practice the second sample after the
event already brings the interval close to its final value.

\end{document}
